\documentclass[11pt]{article}
\usepackage[a4paper,margin=1in]{geometry}
\usepackage{amsmath,amssymb,amsthm}
\usepackage{booktabs}
\usepackage{microtype}
\usepackage[hidelinks]{hyperref}
\usepackage{enumitem}
\usepackage{xcolor}

\newtheorem{proposition}{Proposition}
\newtheorem{lemma}{Lemma}
\newtheorem{remark}{Remark}
\newcommand{\bpos}{\mathrm{bpos}}
\newcommand{\capa}{\mathrm{cap}}
\newcommand{\pen}{\mathrm{pen}}
\newcommand{\ad}{\bar d}
\newcommand{\D}{\mathcal{D}}
\newcommand{\T}{^{\top}}
\DeclareMathOperator{\conv}{conv}

\title{Facility Location with Decision-Dependent Scenario Probabilities:\\
Linked Lifted Cuts and a Scenario-Decomposed Branch-and-Cut}
\author{Karyamsetty Martin}
\date{31 October 2026}

\begin{document}
\maketitle

\begin{abstract}
We study two-stage facility location in which the probability of each demand scenario is an affine function of the
opening decisions, so that the objective $f\T y + \sum_s p_s(y)\,Q_s(y)$ is bilinear in the binary first stage $y$ and the
recourse values $Q_s(y)$. The standard treatment linearizes the products $y_j\theta_s$ by McCormick envelopes over
global bounds on $\theta_s$; on our instances its LP relaxation has a root gap of 6--10\%. We propose a
\emph{linked lifted-cut master}: the posiform of $p_s(y)$ splits each product into factors $y_jQ_s(y)$ and
$(1-y_j)Q_s(y)$, each factor is the value of the recourse LP at a scaled right-hand side, and every vertex of the
shared recourse dual polyhedron therefore yields a cut that is linear in $(y, Y)$ with $Y = yy\T$ lifted. Linking rows
tie the factor epigraphs of one scenario together. The relaxation is exact on binary points and has root gaps of
0.07--0.43\% on the same instances. Because the linked LP is dual block-angular with closed-form scenario blocks, its
value is computed by a stabilized multicut cutting-plane master over $(y,Y,\theta)$ whose cuts are valid at every node
of a branch-and-bound tree. The resulting branch-and-cut proves optimality on all 25 test instances with 25 to 50
sites and 20 scenarios, in 4 s to 12 min, while a McCormick branch-and-cut sharing the same tree code leaves seven of
fifteen unproven after two to four hours and, with node-local recourse bounds, still leaves six of 25 unproven and
needs three to four hundred times more nodes; a full MIP solver closes the McCormick gap faster than that tree but
is still $1.6$--$13\times$ slower than the decomposed solver. A node-parallel MPI version of the
solver reaches $1.6$--$3.3\times$ on seven workers, and a sparse-influence variant of the lifting proves instances
with 500 sites in under four minutes.
\end{abstract}

\section{Introduction}

In the classical two-stage stochastic facility location problem the scenario probabilities are data. In many
applications they are not: opening a flagship store shifts demand toward it, a new depot changes which disruption
scenarios matter, a hub makes the high-volume scenario more likely. Stochastic programs in which the decision
influences the probability measure are called problems with \emph{endogenous uncertainty} of the
decision-dependent-probability type \cite{dupacova,hellemo}. They are hard for a structural reason: the expectation
$\sum_s p_s(y) Q_s(y)$ is a sum of products of a function of $y$ and a function of $y$, and convexity of $Q_s$ is
not inherited by the product.

The published computational approaches either hand the bilinear objective to a global solver \cite{hellemo} or
linearize the products $y_j\theta_s$ of a binary and a bounded continuous variable with McCormick envelopes, which
are exact at binary points \cite{basciftci}. The McCormick relaxation is, however, weak: between a binary point and
the next, the product may sit anywhere in a box whose width is the whole range of the recourse value. The
branch-and-bound tree then has to do what the relaxation does not.

This paper makes three contributions.
\begin{enumerate}[nosep]
\item A \emph{linked lifted-cut} LP relaxation of the problem that is exact on binary points and whose root gap is
      one to two orders of magnitude smaller than the McCormick relaxation on the same instances
      (Section~\ref{sec:linked}).
\item A \emph{scenario decomposition} of that LP in which every scenario block has a closed-form value and
      subgradient, so the node bound of a branch-and-bound tree is a small stabilized cutting-plane master whose
      cuts remain valid throughout the tree (Section~\ref{sec:decomp}).
\item A branch-and-cut built on it, with a computational study against the McCormick master in the same tree code,
      with static and with node-local bounds, and inside a full MIP solver; an ablation of the linking rows;
      sensitivity in the number of scenarios, the strength of the endogeneity and the number of customers, and a
      $b_s = 0$ variant on which the linking rows are vacuous; a C++ implementation; a
      node-parallel MPI version; and a sparse-influence variant of the lifting that reaches 500 sites
      (Sections~\ref{sec:bnc}--\ref{sec:scaling}).
\end{enumerate}

\section{Problem}\label{sec:problem}

Sites $i = 1,\dots,n$ with capacity $\capa_i$ and fixed cost $f_i$ serve customers $k = 1,\dots,m$. Scenario
$s = 1,\dots,S$ has demands $d_s \in \mathbb{R}^m_{+}$ and unit transport costs $C_{ik}$; unmet demand costs $\pen$ per
unit. The first stage is $y \in \{0,1\}^n$. The recourse value is the transportation LP
\begin{equation}\label{eq:Q}
Q_s(y) = \min_{x,u\ge 0}\Big\{ \sum_{ik} C_{ik}x_{ik} + \pen\sum_k u_k \;:\; \sum_i x_{ik} + u_k \ge d_{sk}\ \forall k,\
\sum_k x_{ik} \le \capa_i y_i\ \forall i \Big\},
\end{equation}
whose dual is
\begin{equation}\label{eq:dual}
Q_s(y) = \max_{(\alpha,\beta)\in\D}\ \alpha\T d_s - \sum_i \beta_i \capa_i y_i,\quad
\D = \{(\alpha,\beta): \alpha_k - \beta_i \le C_{ik},\ 0 \le \alpha \le \pen,\ 0 \le \beta \le \pen\}.
\end{equation}
(Demand satisfaction is written as an inequality; since all costs are positive it is equivalent to the equality
form and gives $\alpha\ge0$, and the unmet-demand variables give $\alpha\le\pen$. No primal variable imposes
$\beta\le\pen$: it is a valid truncation. Given any dual-feasible $(\alpha,\beta)$, replacing $\beta_i$ by
$\max(0,\max_k(\alpha_k - C_{ik}))\le\pen$ keeps feasibility, since $C\ge0$, and does not lower the objective, whose
coefficient $-\capa_iy_i$ of $\beta_i$ is nonpositive. Hence the truncation changes no value $Q_s(y)$, $y\ge0$, and
makes $\D$ a polytope even where $y_i = 0$ leaves $\beta_i$ free; $\D$ is the convex hull of its vertices. The rays
of the untruncated dual cone, $\beta_i\to\infty$ where $y_i = 0$, are therefore never needed: for every
$y\in[0,1]^n$ the value $Q_s(y)$ is attained at a vertex of the truncated $\D$, and every cut below is built from
those vertices alone.) Two facts about \eqref{eq:dual} carry the
paper. First, $\D$ does not depend on $y$, so every vertex $(\alpha,\beta)$ of $\D$ gives the Benders cut $Q_s(y) \ge \alpha\T d_s - (\beta\circ\capa)\T y$, valid for every $y$. Second, write
$Q(r, c)$ for the value of \eqref{eq:Q} with demand $r$ and capacity $c$; then $Q$ is positively homogeneous,
$Q(\lambda r, \lambda c) = \lambda Q(r, c)$ for $\lambda \ge 0$, and the same vertices of $\D$ are optimal.

\begin{remark}[scope of the two facts]\label{rem:scope}
Neither fact is specific to transportation. For any second-stage LP $Q(h) = \min\{q\T x : Wx\ge h,\ x\ge0\}$
with fixed recourse matrix $W$ and right-hand side $h(y) = h_0 + Ty$ affine in $y$, the dual feasible set
$\{\pi\ge0 : W\T\pi\le q\}$ does not depend on $y$ and $Q$ is positively homogeneous in $h$, so
$y_jQ(h(y)) = Q(y_jh_0 + T(y_jy))$ and $(1-y_j)Q(h(y)) = Q((1-y_j)h_0 + T(y - y_jy))$ are values of the same LP at
right-hand sides affine in $(y, yy\T)$, and every dual vertex yields a cut linear in the lifted variables. Everything
in Sections~\ref{sec:linked}--\ref{sec:decomp} uses only this and a dual feasible set that is a polytope, or can be
replaced by one without changing the value at any right-hand side that occurs. The first is a structural condition:
the dual of $Q(h)$ is bounded iff $\{x\ge0 : Wx > 0\}$ is nonempty (Motzkin's transposition theorem), which complete
recourse with a penalised slack on every row guarantees. The transportation LP has slacks on the demand rows only and
its capacity duals are unbounded where $y_i = 0$; the truncation $\beta\le\pen$ is valid there because of the
bipartite structure $\alpha_k - \beta_i\le C_{ik}$ with $C\ge0$. Another recourse structure without a slack on every
row needs a truncation argument of its own, or the feasibility rays carried as cuts.
\end{remark}

The scenario probabilities depend on $y$ through an affine map that is a probability vector on every binary point:
\begin{equation}\label{eq:p}
p_s(y) = b_s + a_s\T y,\qquad \sum_s b_s = 1,\quad \sum_s a_s = 0,\quad
\bpos_s := b_s + \sum_j \min(a_{sj},0) \ge 0 .
\end{equation}
The last condition is equivalent to $p_s(y)\ge 0$ for all $y\in\{0,1\}^n$. The problem is
\begin{equation}\label{eq:F}
\min_{y\in\{0,1\}^n}\ F(y) = f\T y + \sum_s p_s(y)\,Q_s(y).
\end{equation}
With exogenous probabilities ($a_s = 0$) this is a standard Benders decomposition. The products $p_s(y)Q_s(y)$ are the
entire difficulty.

\begin{remark}[scenario-dependent recourse data]
Nothing below uses that $C$, $\capa$ and $\pen$ are the same in every scenario. With scenario data $C^s$, $\capa^s$,
$\pen^s$ the dual polyhedron becomes $\D_s$, still independent of $y$, and every statement and proof goes through with
$\D_s$, $\capa^s$ in place of $\D$, $\capa$; the implementation already keeps one vertex list per scenario
(Section~\ref{sec:linked}), so it covers this case as written. We keep the shared polyhedron in the notation only
because it lets one vertex serve every scenario.
\end{remark}

\paragraph{Instances.} All computations use one generator: $m = 15$, $S = 20$, costs $C_{ik}\sim U(1,5)$, demands
$d_{sk}\sim U(1.5,4.5)$ scaled by a factor rising from $0.5$ to $1.6$ across scenarios, total capacity about
$1.8\times$ the mean total demand, $\pen = 10$, and fixed costs calibrated as the recourse saving of each site at the
mean demand times $U(0.6,1.6)$, so that opening a strict subset is optimal. In the ``strong'' endogeneity used
throughout, the low-demand half of the scenarios carries the probability mass at $y = 0$ and each opened site moves a
share of up to 90\% of that mass to the high-demand scenarios: facilities attract demand. Three seeds per size.

\section{The McCormick master and why it is weak}\label{sec:mcc}

Introduce $\theta_s \ge Q_s(y)$ with Benders cuts and expand
$p_s(y)\theta_s = b_s\theta_s + \sum_j a_{sj}\,y_j\theta_s$. With valid bounds $Q_s^{\min}\le\theta_s\le Q_s^{\max}$,
each product $y_j\theta_s$ is replaced by $z_{sj}$ in its McCormick envelope
\begin{equation}\label{eq:mcc}
z_{sj} \ge Q_s^{\min}y_j,\quad z_{sj}\ge \theta_s - Q_s^{\max}(1-y_j),\quad
z_{sj}\le Q_s^{\max}y_j,\quad z_{sj}\le \theta_s - Q_s^{\min}(1-y_j).
\end{equation}
Because $y_j$ is binary the envelope is exact on integer points \cite{basciftci}, so minimizing
$f\T y + \sum_s (b_s\theta_s + \sum_j a_{sj}z_{sj})$ over \eqref{eq:mcc} and the Benders cuts is a correct MILP. We
call it \textsc{MC-B}. The only valid bounds available before solving are $Q_s^{\max} = \pen\sum_k d_{sk} = Q_s(\mathbf 0)$
(nothing open) and $Q_s^{\min} = Q_s(\mathbf 1)$ (everything open); both hold for every $y\in[0,1]^n$ because
$Q_s$ is nonincreasing in $y$ ($\beta\ge0$ in \eqref{eq:dual}). At a fractional $y$ the envelope lets $z_{sj}$ sit anywhere
in a box of width $Q_s^{\max}-Q_s^{\min}$, the full range of the recourse, and the LP picks the cheapest corner.
Nothing in \eqref{eq:mcc} connects the value of $\theta_s$ at a point to the value of the product at that point. On
our instances the root gap of \textsc{MC-B} is 6.1--9.7\% (Table~\ref{tab:main}).

\section{The linked lifted-cut master}\label{sec:linked}

The idea is to give each product $y_jQ_s(y)$ its own exact convex under-estimator rather than relax $y_j\theta_s$
through global bounds. Three observations make this possible.

\paragraph{Posiform.} Since $\bpos_s\ge0$, the probability splits into nonnegative pieces by the sign of $a_{sj}$:
\begin{equation}\label{eq:posiform}
p_s(y) = \bpos_s + \sum_{j: a_{sj}>0} a_{sj}\,y_j + \sum_{j:a_{sj}<0} |a_{sj}|\,(1-y_j)\qquad (y\in\{0,1\}^n),
\end{equation}
so $p_s(y)Q_s(y)$ is a nonnegative combination of three kinds of factors: $Q_s(y)$, $y_jQ_s(y)$ and $(1-y_j)Q_s(y)$.

\paragraph{Factors are recourse LPs.} For $y\in\{0,1\}^n$ let $Y = yy\T$ and let $Y_j$ denote its $j$-th row, so
$Y_{ji} = y_jy_i$ and $Y_{jj} = y_j$. By homogeneity,
\begin{equation}\label{eq:factors}
y_j\,Q_s(y) = Q\big(y_j d_s,\ \capa\circ Y_j\big),\qquad (1-y_j)\,Q_s(y) = Q\big((1-y_j)d_s,\ \capa\circ(y - Y_j)\big).
\end{equation}
Both are values of the same dual polyhedron $\D$ at a different right-hand side. Hence every vertex $(\alpha,\beta)$
of $\D$, with $g := \beta\circ\capa$, yields cuts that are \emph{linear in $(y,Y)$}:
\begin{align}
e^0_s &\ge \alpha\T d_s - g\T y, \label{eq:cut0}\\
e^+_{sj} &\ge \alpha\T d_s\, y_j - g\T Y_j, \label{eq:cutp}\\
e^-_{sj} &\ge \alpha\T d_s\,(1-y_j) - g\T(y - Y_j). \label{eq:cutm}
\end{align}
\begin{remark}[conditional configurations]\label{rem:cond}
At a master point $(\hat y,\hat Y)$ with $0<\hat y_j<1$, homogeneity gives
$Q(\hat y_j d_s, \capa\circ\hat Y_j) = \hat y_j\,Q_s(\hat Y_j/\hat y_j)$ and
$Q((1-\hat y_j)d_s, \capa\circ(\hat y-\hat Y_j)) = (1-\hat y_j)\,Q_s\big((\hat y-\hat Y_j)/(1-\hat y_j)\big)$, and the
McCormick rows $0\le Y_{jl}\le\min(y_j,y_l)$, $y_l - Y_{jl}\le 1-y_j$ put both quotients in $[0,1]^n$. The factor
cuts are therefore ordinary Benders cuts evaluated at the \emph{conditional configurations} $\hat Y_j/\hat y_j$ and
$(\hat y-\hat Y_j)/(1-\hat y_j)$, the points the master believes in given $y_j = 1$ and given $y_j = 0$, exactly the
points of Balas's lift-and-project and of level-1 RLT. This is why the vertices separated at the factor points are
so effective and why the root pool is nearly complete for the tree (Section~\ref{sec:bnc}).
\end{remark}

A vertex found for one factor of one scenario is a valid cut for every factor of every scenario. The implementation
keeps one list of vertices per scenario, $\{(\alpha_k,\beta_k)\}_{k\le K_s}$, each list evaluated against its own
$d_s$; a vertex separated for scenario $s$ is added to the list of $s$ only, which is a choice of economy rather than
of validity (adding it to every list is sound and was not needed).

\paragraph{Linking.} For every $j$, $y_jQ_s(y) + (1-y_j)Q_s(y) = Q_s(y)$. The master carries this as the rows
\begin{equation}\label{eq:link}
e^0_s = e^+_{sj} + e^-_{sj}\qquad \forall s,\ \forall j.
\end{equation}
Without them each factor epigraph is bounded only by its own cuts and the three epigraphs of a scenario can sit at
inconsistent levels; with them the relaxation tightens substantially (we call the model without \eqref{eq:link}
\emph{unlinked}; Section~\ref{sec:results} reports its root gap as an ablation). Proposition~\ref{prop:block} below
makes precise how much they tighten: the effect of the linking rows on the value of scenario $s$ is
$\bpos_s$ times the amount by which they lift $e^0_s$, so a scenario with $\bpos_s = 0$ gains nothing from them.

\paragraph{The master LP.} With pair variables $Y_{jl} = Y_{lj}$, $j<l$, and $Y_{jj} := y_j$, the linked lifted-cut
master \textsc{C2L} is
\begin{equation}\label{eq:c2l}
\begin{aligned}
\min\ & f\T y + \sum_s \Big[\bpos_s\, e^0_s + \sum_{a_{sj}>0} a_{sj}\,e^+_{sj} + \sum_{a_{sj}<0}|a_{sj}|\,e^-_{sj}\Big]\\
\text{s.t.}\ & Y_{jl}\le y_j,\quad Y_{jl}\le y_l,\quad Y_{jl}\ge y_j+y_l-1,\quad Y_{jl}\ge 0 \qquad (j<l),\\
& \text{\eqref{eq:cut0}--\eqref{eq:cutm} for every vertex in the pool},\quad \eqref{eq:link},\quad
y\in[0,1]^n,\ e\ge 0 .
\end{aligned}
\end{equation}

\begin{proposition}[validity and exactness]\label{prop:exact}
Let $v(y)$ be the value of \eqref{eq:c2l} with $y$ fixed. Then $v(y)\le F(y)$ for every $y\in[0,1]^n$, and $v(\hat y) = F(\hat y)$ for every $\hat y\in\{0,1\}^n$ at which the pool contains, for each
$s$, a vertex optimal in \eqref{eq:dual} at $\hat y$.
\end{proposition}
\begin{proof}
Validity: fix $y\in[0,1]^n$ and set $Y_{jl} = y_jy_l$, which satisfies the McCormick rows for every $y$ in the
box. Set $e^0_s = Q_s(y)$, $e^+_{sj} = y_jQ_s(y)$ and $e^-_{sj} = (1-y_j)Q_s(y)$. By positive homogeneity,
$y_jQ_s(y) = Q(y_jd_s,\ y_j\,\capa\circ y) = Q(y_jd_s,\ \capa\circ Y_j)$ and likewise for $(1-y_j)Q_s(y)$, so each of
\eqref{eq:cut0}--\eqref{eq:cutm} is the Benders cut of the corresponding LP at its own optimum and holds; the
linking rows \eqref{eq:link} hold with equality. The objective at this point is $F(y)$ by \eqref{eq:posiform} (which
is an identity for every $y\in[0,1]^n$, not only binary ones, because $\sum_j\min(a_{sj},0)$ and the two sums
recombine to $b_s + a_s\T y$). Hence $v(y)\le F(y)$. Exactness: at $\hat y$ with the optimal vertex $(\alpha^*,\beta^*)$ in the pool,
cut \eqref{eq:cut0} forces $e^0_s\ge Q_s(\hat y)$. If $\hat y_j = 1$ then $Y_j = \hat y$ and \eqref{eq:cutp} at
$(\alpha^*,\beta^*)$ forces $e^+_{sj}\ge Q_s(\hat y)$; if $\hat y_j = 0$ then $\hat y - Y_j = \hat y$ and
\eqref{eq:cutm} forces $e^-_{sj}\ge Q_s(\hat y)$. These lower bounds are attainable together with
\eqref{eq:link} (set the other factor of each pair to zero), every objective coefficient is nonnegative, so the
minimum is attained at them and equals $F(\hat y)$ by \eqref{eq:posiform}.
\end{proof}

\begin{remark}
\eqref{eq:cutp}--\eqref{eq:cutm} are what one obtains by multiplying the dual constraint system by $y_j$ and
$(1-y_j)$ and linearizing, i.e.\ the mixed-binary reformulation--linearization of Adams and Sherali
\cite{adams-sherali} applied on the \emph{dual} side of the recourse. The dual side matters: the polyhedron $\D$ is
shared, so the lifted model is generated by cuts from one pool instead of carrying $S(1+2n)$ copies of the primal
recourse, and, as the next section shows, its scenario blocks have closed-form values.
\end{remark}

\paragraph{Dominance over the McCormick master.} Let \textsc{C2R} denote \eqref{eq:c2l} without the linking rows
\eqref{eq:link} (the \emph{unlinked} model). Write $v_{\textsc{MC}}(y)$, $v_{\textsc{C2R}}(y)$ and $v_{\textsc{C2L}}(y)$ for the
values of the three masters with $y$ fixed, all three built from the same set of dual vertices.

\begin{proposition}[dominance]\label{prop:dom}
If the vertex set contains the zero vertex and a vertex optimal in \eqref{eq:dual} at $y = \mathbf 1$, then for
every $y\in[0,1]^n$
\[
v_{\textsc{MC}}(y)\ \le\ v_{\textsc{C2R}}(y)\ \le\ v_{\textsc{C2L}}(y)\ \le\ F(y),
\]
and the same chain holds for the LP relaxations (minima over $y\in[0,1]^n$).
\end{proposition}
\begin{proof}
The last inequality is Proposition~\ref{prop:exact}; the middle one holds because \textsc{C2L} adds constraints
to \textsc{C2R}. For the first, fix $y$ and let $(Y,e)$ be optimal for \textsc{C2R} at $y$; by
Proposition~\ref{prop:block} without the linking step, $e^0_s = L_0$, $e^+_{sj} = L^+_j$, $e^-_{sj} = L^-_j$ with the
levels \eqref{eq:levels} at $(y,Y)$. Let $(\alpha^1,\beta^1)$ be the vertex optimal at $\mathbf 1$, so
$\ad^{\,1} - g^{1\top}\mathbf 1 = Q_s^{\min}$. Define a point of \textsc{MC-B}: $\theta_s = L_0$, and for each
$j$ with $a_{sj}>0$, $z_{sj} = \max\{Q_s^{\min}y_j,\ \theta_s - Q_s^{\max}(1-y_j)\}$; for $a_{sj}<0$,
$z_{sj} = \min\{Q_s^{\max}y_j,\ \theta_s - Q_s^{\min}(1-y_j)\}$.

\emph{Feasibility.} $\theta_s$ satisfies every Benders cut since $L_0$ is their maximum. $\theta_s\ge
\ad^{\,1} - g^{1\top}y\ge \ad^{\,1} - g^{1\top}\mathbf 1 = Q_s^{\min}$ because $g^1\ge0$ and $y\le\mathbf 1$;
$\theta_s\le\max_k \ad_k\le Q_s(\mathbf 0) = Q_s^{\max}$ by weak duality at $y = \mathbf 0$. With
$\theta_s\in[Q_s^{\min},Q_s^{\max}]$ the envelope \eqref{eq:mcc} is nonempty and $z_{sj}$ lies in it by
construction.

\emph{Objective.} Using $b_s = \bpos_s + \sum_{a_{sj}<0}|a_{sj}|$, the block objective of \textsc{MC-B} at this
point is $\bpos_s L_0 + \sum_{a_{sj}>0}a_{sj}z_{sj} + \sum_{a_{sj}<0}|a_{sj}|(L_0 - z_{sj})$, so it suffices to show
$z_{sj}\le L^+_j$ when $a_{sj}>0$ and $L_0 - z_{sj}\le L^-_j$ when $a_{sj}<0$. Let $k$ be the pool index attaining $L_0$ (the zero vertex when $L_0 = 0$, so that $\ad_k = 0$ and
$g_k = 0$) and recall the McCormick rows give $0\le Y_{jl}\le \min(y_j,y_l)$ and $y_l - Y_{jl}\le 1-y_j$.
\begin{itemize}[nosep]
\item $L^+_j\ge \ad^{\,1}y_j - g^{1\top}Y_j\ge y_j(\ad^{\,1} - g^{1\top}\mathbf 1) = Q_s^{\min}y_j$.
\item $L_0 - L^+_j\le (\ad_k - g_k\T y) - (\ad_ky_j - g_k\T Y_j) = \ad_k(1-y_j) -
      g_k\T(y - Y_j)\le \ad_k(1-y_j)\le Q_s^{\max}(1-y_j)$.
\item $L^-_j\ge \ad^{\,1}(1-y_j) - g^{1\top}(y-Y_j)\ge (1-y_j)(\ad^{\,1} - g^{1\top}\mathbf 1) = Q_s^{\min}(1-y_j)$.
\item $L_0 - L^-_j\le (\ad_k - g_k\T y) - (\ad_k(1-y_j) - g_k\T(y-Y_j)) = \ad_ky_j -
      g_k\T Y_j\le Q_s^{\max}y_j$.
\end{itemize}
The first two give $z_{sj}\le L^+_j$ for $a_{sj}>0$; the last two give $z_{sj}\ge L_0 - L^-_j$ for $a_{sj}<0$.
Hence the \textsc{MC-B} objective at the constructed point is at most $v_{\textsc{C2R}}(y)$, and
$v_{\textsc{MC}}(y)$, a minimum, is smaller still.
\end{proof}

The proof shows where the strength comes from: the lifted cuts bound each product $y_jQ_s(y)$ by the recourse
value at a scaled right-hand side, which is at least as tight as any corner of the envelope \eqref{eq:mcc} at every
point and strictly tighter whenever the active dual vertex of the factor differs from the one active at
$\mathbf 0$ or $\mathbf 1$. Section~\ref{sec:results} quantifies the two steps: on our instances the unlinked model
already reduces the root gap from 6--10\% to 0.26--1.14\%, and the linking rows reduce it a further
$1.7$--$6\times$.

\paragraph{Sparse influence.} Let $\mathcal I = \{j : a_{sj}\ne 0 \text{ for some } s\}$ be the influencing sites.
A factor $y_jQ_s(y)$ or $(1-y_j)Q_s(y)$ with $j\notin\mathcal I$ has zero weight in the objective of \eqref{eq:c2l}
and enters only through its linking row \eqref{eq:link}, which raises $e^0_s$ to at least $e^+_{sj}+e^-_{sj}$. The
\emph{sparse} master drops the factors of $j\notin\mathcal I$ altogether: cuts \eqref{eq:cutp}--\eqref{eq:cutm} and
rows \eqref{eq:link} are kept for $j\in\mathcal I$ only. It is a relaxation of \eqref{eq:c2l}, hence valid, and
Proposition~\ref{prop:exact} still holds because its exactness argument uses the $e^0$ cuts and the factors of
influencing sites only. Its pair columns are $\{(j,l): j\in\mathcal I \text{ or } l\in\mathcal I\}$, that is
$|\mathcal I|\,n - |\mathcal I|(|\mathcal I|+1)/2$ instead of $n(n-1)/2$, and the scenario blocks of
Section~\ref{sec:decomp} evaluate only the rows $Y_j$, $j\in\mathcal I$, of the lifted matrix. When every site
influences demand the two masters coincide.

On our instances the root gap of \eqref{eq:c2l} is 0.07--0.43\% against 6.1--9.7\% for \textsc{MC-B}
(Table~\ref{tab:main}). The price is size: $n(n-1)/2$ pair columns and cut rows that are dense over them (about
720 nonzeros per row at $n = 50$). Re-solving the monolithic LP at every node is already slow at 35--40 sites;
the decomposition removes that cost.

\section{Scenario decomposition of the linked master}\label{sec:decomp}

Fix the coupling variables $(y,Y)$. Scenario $s$ owns the private variables $(e^0_s, e^+_s, e^-_s)$, its cuts and
its linking rows, and nothing else, so \eqref{eq:c2l} is dual block-angular. With a pool of $K$ vertices for scenario
$s$, write $\ad_k = \alpha_k\T d_s$ and $g_k = \beta_k\circ\capa$. The tightest epigraph levels are
\begin{equation}\label{eq:levels}
\begin{aligned}
L_0 &= \max\Big(0,\ \max_k\, \ad_k - g_k\T y\Big),\\
L^+_j &= \max\Big(0,\ \max_k\, \ad_k\, y_j - g_k\T Y_j\Big),\\
L^-_j &= \max\Big(0,\ \max_k\, \ad_k\,(1-y_j) - g_k\T(y-Y_j)\Big).
\end{aligned}
\end{equation}
The $0$ in each maximum is the cut of the zero vertex $(\alpha,\beta) = (0,0)\in\D$, feasible for every instance and
kept in every list, so each level is attained by a vertex of the pool.

\begin{proposition}[closed-form block]\label{prop:block}
For fixed $(y,Y)$ the value of the block of scenario $s$ in \eqref{eq:c2l} is
\begin{equation}\label{eq:phi}
\varphi_s(y,Y) = \bpos_s\, e^{0*}_s + \sum_{a_{sj}>0} a_{sj}\,L^+_j + \sum_{a_{sj}<0}|a_{sj}|\,L^-_j,\qquad
e^{0*}_s = \max\Big(L_0,\ \max_j\, L^+_j + L^-_j\Big).
\end{equation}
\end{proposition}
\begin{proof}
All objective coefficients of the block are nonnegative, so each variable is pushed to its lower bound subject to
\eqref{eq:link}. Row $j$ of \eqref{eq:link} is feasible with $e^+_{sj}\ge L^+_j$, $e^-_{sj}\ge L^-_j$ iff
$e^0_s\ge L^+_j + L^-_j$, which gives $e^{0*}_s$. Given $e^0_s = e^{0*}_s$: if $a_{sj}>0$ set $e^+_{sj} = L^+_j$ and
$e^-_{sj} = e^{0*}_s - L^+_j \ge L^-_j$, which costs $a_{sj}L^+_j$ and is the least possible; if $a_{sj}<0$ the
mirror choice costs $|a_{sj}|L^-_j$; if $a_{sj}=0$ the row is free.
\end{proof}

\begin{remark}[what the linking rows buy]\label{rem:bpos}
For a pair $(s,j)$ at most one of $e^+_{sj}$, $e^-_{sj}$ carries a nonzero objective coefficient; the other appears
only in its own cuts and in row $j$ of \eqref{eq:link}. Eliminating it, $e^-_{sj} = e^0_s - e^+_{sj}\ge L^-_j$ when
$a_{sj}>0$, shows that the linking rows do nothing but impose $e^0_s\ge\max_j(L^+_j+L^-_j)$: the active factor
epigraphs stay at their own levels, and the block value differs from the unlinked one, $\bpos_sL_0 + \sum a_{sj}L^+_j
+ \sum|a_{sj}|L^-_j$, exactly by $\bpos_s\,(e^{0*}_s - L_0)$. The posiform objective of \eqref{eq:c2l} is what one
gets by substituting $e^+_{sj} = e^0_s - e^-_{sj}$ for $a_{sj}<0$ into $b_se^0_s + \sum_j a_{sj}e^+_{sj}$, and
$\bpos_s$ is the weight that remains on $e^0_s$ after that substitution. Hence the tightening from \textsc{C2R} to
\textsc{C2L} is carried entirely by the residual base mass: a scenario with $\bpos_s = 0$, which happens when every
unit of its base probability can be moved away by closing sites, is bounded as in the unlinked model. In the generator
of Section~\ref{sec:problem} every scenario has $\bpos_s\approx 0.009$: a tenth of each low-demand scenario's base
mass $0.091$ is immovable, and each high-demand scenario keeps its whole base mass $0.009$ since nothing is moved
away from it. The weight on the linking is therefore small; that it still removes two thirds of the unlinked gap (Section~\ref{sec:results}) is
because $e^{0*}_s - L_0$ is large at fractional points, where the factor levels $L^+_j + L^-_j$ sit well above $L_0$.
\end{remark}

Each level in \eqref{eq:levels} is a maximum of affine functions of $(y,Y)$, so $\varphi_s$ is convex
piecewise-linear in $(y,Y)$, and a subgradient is read off the active pieces (envelope theorem). Let $k_0$, $k^+_j$,
$k^-_j$ be the indices attaining the maxima in \eqref{eq:levels} (a level equal to $0$ has gradient $0$), and let
$j^*$ attain $\max_j L^+_j+L^-_j$. Ties, among the $j$ or between $L_0$ and $L^+_{j^*}+L^-_{j^*}$, are exactly the
points where $\varphi_s$ is not differentiable; any choice of the active piece gives an element of the
subdifferential of a pointwise maximum of convex functions, and the master cut \eqref{eq:mcut} is valid for any element
of it. With $Y_{jj} = y_j$ and $Y_{jl}$ stored once for $j<l$, the gradients of the three
levels are, writing $\partial_{y_l}$ and $\partial_{Y_{jl}}$ for the coordinates,
\begin{equation}\label{eq:grad}
\begin{aligned}
\nabla L_0:\ & \partial_{y_l} = -g_{k_0,l}\quad(\text{all } l);\\
\nabla L^+_j:\ & \partial_{y_j} = \ad_{k^+_j} - g_{k^+_j,j},\qquad \partial_{Y_{jl}} = -g_{k^+_j,l}\quad (l\ne j);\\
\nabla L^-_j:\ & \partial_{y_j} = -\ad_{k^-_j},\qquad \partial_{y_l} = -g_{k^-_j,l}\quad (l\ne j),\qquad \partial_{Y_{jl}} = +g_{k^-_j,l}\quad (l\ne j),
\end{aligned}
\end{equation}
where $\partial_{Y_{jl}}$ is the derivative with respect to the stored pair $Y_{\min(j,l)\max(j,l)}$, contributions of
rows $j$ and $l$ to the same stored pair add, and the diagonal $Y_{jj} = y_j$ has been folded into $\partial_{y_j}$:
in $L^+_j$ the term $-g_{k,j}Y_{jj}$ contributes $-g_{k,j}$ to $\partial_{y_j}$, in $L^-_j$ the term
$-g_{k,j}(y_j - Y_{jj})$ vanishes. The subgradient of $\varphi_s$ is then
\begin{equation}\label{eq:subgrad}
\gamma_s = \bpos_s\,\nabla e^{0*}_s + \sum_{a_{sj}>0}a_{sj}\nabla L^+_j + \sum_{a_{sj}<0}|a_{sj}|\nabla L^-_j,\qquad
\nabla e^{0*}_s = \begin{cases}\nabla L_0 & \text{if } e^{0*}_s = L_0,\\ \nabla L^+_{j^*}+\nabla L^-_{j^*} & \text{otherwise.}\end{cases}
\end{equation}
Evaluating all $S$ blocks costs one matrix product $G_s Y\T$ per scenario and nothing else; no LP is solved.

\paragraph{Node bound.} The value of \eqref{eq:c2l} at a node with box $[\ell,u]$ on $y$ is
\begin{equation}\label{eq:node}
\min\Big\{ f\T y + \sum_s \varphi_s(y,Y)\ :\ \ell\le y\le u,\ (y,Y)\ \text{in the McCormick set}\Big\},
\end{equation}
a convex problem which we solve as a multicut cutting-plane (Benders) master over $(y,Y,\theta)$, $\theta\in\mathbb{R}^S$, with
the McCormick rows and the cuts
\begin{equation}\label{eq:mcut}
\theta_s\ \ge\ \varphi_s(\hat y,\hat Y) + \gamma_s\T\big((y,Y)-(\hat y,\hat Y)\big),\qquad \gamma_s\ \text{from \eqref{eq:subgrad} at } (\hat y,\hat Y).
\end{equation}
Its value is a lower bound on \eqref{eq:node} at every iteration; the best $f\T z + \sum_s\varphi_s(z)$
seen is an upper bound; they meet at the optimum. When the loop stops at tolerance $10^{-6}$ it is the master LP
value, the dual-valid side, that is recorded as the node's bound in the tree, never the oracle value. The LP solver only ever sees this master, which has
$n + n(n-1)/2 + S$ columns, $3n(n-1)/2$ McCormick rows and a few thousand active cuts.

\paragraph{Stabilization and engineering.} Plain Kelley steps oscillate on this master. We generate cuts at the
in-out point $z = \lambda x + (1-\lambda)c$ \cite{benameur}, where $x$ is the master solution, $c$ the center (best
point so far) and $\lambda = 0.2$; if the cut at $z$ does not cut off $x$ a Kelley step at $x$ follows. The center
persists across nodes of the tree, projected onto the new box and the McCormick set, so each node starts from a point
already near-optimal for its parent. Cut rows that have been basic and slack for more than five consecutive solves
are deleted, which keeps the master at 1,000--4,000 rows. Each solve is a warm dual simplex from the previous basis
with devex pricing (HiGHS 1.15.1 \cite{highs}); the solver state is reset before every solve, which costs nothing
measurable and avoids a rare stall of warm re-solves.

\begin{remark}
The decomposition computes the same bound as the monolithic LP \eqref{eq:c2l}; it is a different way to compute it.
Its advantage is that the growth of the cut pool never enters the LP the solver sees: new dual vertices only change
the closed-form blocks.
\end{remark}

\section{Branch-and-cut}\label{sec:bnc}

\paragraph{Root.} Seed the pool with the dual vertices at $y = \mathbf 0$ and $y = \mathbf 1$. Alternate: compute
the bound \eqref{eq:node} at the root box to tolerance; separate at its point by solving one recourse LP per factor
($S(1+2n)$ LPs, parallel over scenarios) and adding every violated vertex to the pool; repeat until nothing is
violated. The bound is first computed to a loose tolerance $10^{-4}$ and tightened to $10^{-6}$ only once separation
finds nothing. Rounding the root point gives the first incumbent.

\paragraph{Tree.} Best-first on the parent bound; branching on the most fractional coordinate. At a node the bound runs
from the persistent center with the node box and an \emph{incumbent cutoff}: as soon as the master value exceeds
the incumbent the node is pruned without finishing the bound, since the master value is a valid lower bound at every
iteration.

\paragraph{Cuts are global.} $\varphi_s$ does not depend on the node box, so every scenario cut in the master and
every dual vertex in the pool is valid at every node. Nothing is undone on backtracking.

\paragraph{Integer leaves.} At a binary $\hat y$ the McCormick rows force $Y = \hat y\hat y\T$, and the levels
collapse: $L^+_j = \hat y_jL_0$, $L^-_j = (1-\hat y_j)L_0$, so $\varphi_s(\hat y,\hat y\hat y\T) = p_s(\hat y)\,L_0$.
The leaf is checked with $S$ recourse LPs: if $Q_s(\hat y) > L_0$ for some $s$ the new vertex is added and the node
re-bounded from its center; otherwise $L_0 = Q_s(\hat y)$ for all $s$, the node bound equals $F(\hat y)$ by
Proposition~\ref{prop:exact}, and $\hat y$ updates the incumbent. This is the only place where exactness of the final
answer is established.

\paragraph{Why fractional nodes do not separate.} New dual vertices are separated only at the root (every factor,
$S(1+2n)$ LPs per round) and at integer leaves ($S$ LPs). On the 15 instances of Table~\ref{tab:main} the root
generates 9,200 to 46,100 vertices per instance and the whole tree adds 0 to 19 more; on nine instances the tree
adds at most one. The root pool is in practice complete for the tree, because the lifted cuts of every factor at the
converged root point already contain the vertices that are optimal in the neighbourhood the tree explores, and any
vertex still missing at an integer leaf is caught there before the leaf can become an incumbent. Separating at
fractional nodes would cost $S(1+2n)$ LPs per node for cuts the root has already found.

\paragraph{Heuristic and proof.} Every fractional node point is rounded once per distinct rounded vector and
evaluated exactly with $S$ recourse LPs, whose optimal dual vertices go into the pool at no extra cost, so the pool
holds an optimal vertex for every integer point the heuristic has visited; on our instances the optimum is found within the first 30 nodes and the rest
of the tree is proof. The run is proven when the heap is empty. Under a wall limit the smallest open bound is reported.

The McCormick solver used for comparison is the same tree code with the node bound replaced by the incremental
\textsc{MC-B} LP, the same heuristic and the same tolerance. The two solvers differ only in the relaxation.

\section{Computational results}\label{sec:results}

All runs on one 8-core laptop (Apple silicon, 16 GB), C++ solver with HiGHS 1.15.1 for all LPs, eight threads for the
recourse LPs, tolerance $10^{-6}$, wall limits 7,200 s at 25--40 sites and 14,400 s at 50, idle machine. Three
masters are run through the \emph{same} branch-and-cut code, heuristic and tolerance: \textsc{MC-B} with the static
bounds $[Q_s(\mathbf 1),\,\pen\sum_k d_{sk}]$ of Section~\ref{sec:mcc}; \textsc{MC-B} with \emph{node-local} bounds,
where at every node $\theta_s$ is bounded by $[Q_s(u),\,Q_s(\ell)]$ for the node box $\ell\le y\le u$ (monotonicity of
$Q_s$ in $y$), the four McCormick rows per $(s,j)$ are rewritten with the node constants, and the recourse LPs this
costs are charged to the run (a child changes one side of the box, so one of $Q_s(u)$, $Q_s(\ell)$ is inherited from
the parent and $S = 20$ LPs are solved per node); and the decomposed linked master of Section~\ref{sec:decomp}. The
McCormick masters separate Benders cuts once at every fractional node point and at integer points; the rounding
heuristic and its dual vertices are shared by all three. Five seeds per size; the static master, which proves
nothing above 35 sites, is run on seeds 1--3 only. The unlinked master of Section~\ref{sec:linked} (cuts
\eqref{eq:cutp}--\eqref{eq:cutm} without rows \eqref{eq:link}) is reported at the root only.

\begin{table}[ht]
\centering\small
\caption{Root relaxation gaps, $(F^* - v)/F^*$, of the three masters at the converged root LP, seeds 1--3; the
unlinked master is the decomposed solver with the linking rows dropped (\texttt{--unlinked}).}
\label{tab:root}
\begin{tabular}{rr rrr @{\hspace{2em}} rr rrr}
\toprule
$n$ & seed & \textsc{MC-B} & unlinked & linked & $n$ & seed & \textsc{MC-B} & unlinked & linked\\
\midrule
% TAB1-BEGIN (generated by solver/build/make_tables)
25 & 1 & 8.29\% & 1.14\% & 0.37\% & 35 & 3 & 7.85\% & 0.60\% & 0.23\%\\
25 & 2 & 7.19\% & 0.73\% & 0.14\% & 40 & 1 & 9.66\% & 0.36\% & 0.10\%\\
25 & 3 & 7.86\% & 1.11\% & 0.31\% & 40 & 2 & 8.69\% & 0.40\% & 0.11\%\\
30 & 1 & 6.75\% & 0.58\% & 0.29\% & 40 & 3 & 8.05\% & 0.87\% & 0.20\%\\
30 & 2 & 7.76\% & 0.41\% & 0.07\% & 50 & 1 & 8.47\% & 0.68\% & 0.16\%\\
30 & 3 & 6.71\% & 0.42\% & 0.22\% & 50 & 2 & 8.34\% & 0.26\% & 0.11\%\\
35 & 1 & 6.13\% & 0.49\% & 0.14\% & 50 & 3 & 6.43\% & 0.46\% & 0.16\%\\
35 & 2 & 8.09\% & 0.73\% & 0.43\% & & & & & \\
% TAB1-END
\bottomrule
\end{tabular}
\end{table}

\paragraph{Root gaps.} Table~\ref{tab:root} separates the two ingredients of the linked master. The lifting alone
takes the root gap from 6--10\% to 0.26--1.14\%; the linking rows divide it by a further $1.7$--$6\times$ (mean about
$3\times$), to 0.07--0.43\%. The ordering \textsc{MC-B} $\le$ unlinked $\le$ linked holds on every instance, as
Proposition~\ref{prop:dom} requires. Node-local bounds do not change the root of \textsc{MC-B}: at the root
$\ell = \mathbf 0$, $u = \mathbf 1$ and the two bound choices coincide.

\begin{table}[ht]
\centering\small
\caption{Branch-and-cut: nodes and wall seconds to proven optimality, same C++ tree code for all three masters.
``Unproven'' gives the gap at the wall limit (7,200 s at 25--40 sites, 14,400 s at 50); the static master is not run on
seeds 4--5. For the decomposed solver ``root'' is the time of the root cutting-plane loop, included in ``time''.}
\label{tab:main}
\begin{tabular}{rr rr rr rrr}
\toprule
& & \multicolumn{2}{c}{\textsc{MC-B}, static} & \multicolumn{2}{c}{\textsc{MC-B}, node-local} & \multicolumn{3}{c}{decomposed}\\
\cmidrule(lr){3-4}\cmidrule(lr){5-6}\cmidrule(lr){7-9}
$n$ & seed & nodes & time (s) & nodes & time (s) & nodes & root (s) & time (s)\\
\midrule
% TAB2-BEGIN (generated by solver/build/make_tables)
25 & 1 & 4,859 & 232 & 1,211 & 36 & 71 & 3 & 15\\
25 & 2 & 2,127 & 39 & 485 & 7 & 9 & 2 & 4\\
25 & 3 & 3,187 & 125 & 425 & 8 & 19 & 5 & 11\\
25 & 4 &  &  & 2,263 & 102 & 21 & 6 & 12\\
25 & 5 &  &  & 1,423 & 40 & 129 & 4 & 26\\
30 & 1 & 21,785 & 4,499 & 2,191 & 101 & 53 & 5 & 24\\
30 & 2 & 5,599 & 326 & 663 & 15 & 57 & 10 & 22\\
30 & 3 & 7,087 & 337 & 1,531 & 39 & 27 & 4 & 10\\
30 & 4 &  &  & 2,207 & 113 & 125 & 28 & 82\\
30 & 5 &  &  & 2,389 & 91 & 109 & 8 & 39\\
35 & 1 & 26,146 & unproven (0.34\%) & 2,479 & 125 & 37 & 7 & 36\\
35 & 2 & 37,905 & 6,715 & 3,471 & 188 & 431 & 15 & 249\\
35 & 3 & 13,007 & 1,133 & 1,827 & 66 & 143 & 11 & 91\\
35 & 4 &  &  & 1,395 & 38 & 311 & 8 & 133\\
35 & 5 &  &  & 223 & 5 & 67 & 45 & 80\\
40 & 1 & 15,027 & unproven (2.63\%) & 5,675 & 763 & 21 & 17 & 53\\
40 & 2 & 15,248 & unproven (2.06\%) & 18,620 & unproven (0.15\%) & 91 & 18 & 94\\
40 & 3 & 15,152 & unproven (1.45\%) & 8,303 & 1,185 & 101 & 12 & 155\\
40 & 4 &  &  & 15,408 & unproven (0.41\%) & 117 & 18 & 163\\
40 & 5 &  &  & 13,965 & 2,600 & 317 & 16 & 452\\
50 & 1 & 16,048 & unproven (3.39\%) & 23,233 & 13,863 & 125 & 35 & 394\\
50 & 2 & 21,859 & unproven (2.50\%) & 27,386 & unproven (0.73\%) & 63 & 22 & 193\\
50 & 3 & 19,222 & unproven (3.23\%) & 16,852 & unproven (1.32\%) & 281 & 38 & 706\\
50 & 4 &  &  & 19,995 & unproven (0.84\%) & 113 & 87 & 517\\
50 & 5 &  &  & 23,902 & unproven (0.40\%) & 205 & 34 & 448\\
% TAB2-END
\bottomrule
\end{tabular}
\end{table}

\paragraph{Trees.} The decomposed solver proves all 25 instances in tens to hundreds of nodes. Static
\textsc{MC-B} needs thousands of nodes at 25 sites and tens of thousands from 30; it proves every 30-site instance
(one in 75 minutes), two of three at 35, and none at 40 or 50, where after two and four hours its gaps are
1.5--3.4\% and its best solution is not the optimum on three of the six instances. Node-local bounds are the fair
comparison and they change the picture at small sizes: the McCormick tree shrinks by a factor of $4$--$8$ against
the static tree in the same code at 25 sites, and at 25--30 sites its time is within a factor of ten of the
decomposed solver's on every seed, faster on two of ten. What node-local bounds do not change is the growth rate. The
McCormick tree at 25, 30, 35, 40, 50 sites holds about 700, 1,500, 2,600, 11,000, 22,000 nodes on seeds 1--3
(unproven runs counted at the wall), a factor of about two per five sites, because every node still starts from the
6--10\% root gap and tightens only by branching; the decomposed tree stays in the tens to hundreds because its node
gap is 0.1--0.5\% before the first branch. At 40 sites the node-local master is $6$--$14\times$ slower on the three
seeds it proves and unproven after two hours on seeds 2 and 4 (gaps 0.15\% and 0.41\%); at 50 sites it proves seed
1 alone, in 13,863 s against 394 s, and leaves the other four unproven after four hours at 0.40--1.32\%. It always
holds the optimal solution; what it lacks is the proof. The node counts, not the wall times, are the cleanest statement
of what the relaxation buys. Seeds 4 and 5 fall inside the ranges of seeds 1--3 for the decomposed solver (root gaps
0.08--0.55\%, 12--517 s).

\paragraph{McCormick inside a full MIP solver.} To rule out a weak implementation of the comparison, static
\textsc{MC-B} was also handed to the HiGHS MIP solver: the MILP is built with the complete cut pool of the converged
\textsc{MC-B} root LP of Table~\ref{tab:root}, solved to integer optimality, the integer point is separated with
$S$ recourse LPs, the new cuts are appended and the MILP re-solved (HiGHS has no lazy-constraint callback, so each
round is a fresh MIP solve; 1,800 s per solve, 7,200 s wall). On the six 30- and 40-site instances of seeds 1--3 it
proves all six in three to seven rounds: 180, 36 and 30 s at 30 sites and 234, 1,264 and 584 s at 40, against 24, 22,
10, 53, 94 and 155 s for the decomposed solver in Table~\ref{tab:main}. The MIP solver, with its own cutting planes,
presolve and primal heuristics, closes the McCormick root gap two to twenty-five times faster than our node-local
\textsc{MC-B} tree, and it still needs 1.6 to 13 times the wall time of the decomposed solver on every instance
(2,328 s against 358 s summed), whose tree has none of that machinery. Restarting the MIP every round discards its
tree, which a lazy-constraint callback would keep: on the slowest instance the six MIP solves take 75 to 310 s each,
every one a tree built from scratch at 40 sites. The experiment shows that a mature
MIP code narrows the margin but does not remove it, because the gap it has to close by branching is the same 6.7--9.7\%
of Table~\ref{tab:root} that the plain tree faces.
Seeding the node-local McCormick master with the converged linked root pool instead of its own, as a referee
suggested, changes nothing measurable: on the nine instances with 25 to 35 sites the McCormick root value is the same
to the printed digits, because the converged McCormick root already satisfies every Benders cut at its own point and
the linked pool contains no row the McCormick master can express, and the trees differ from the unseeded ones by at
most 3\% in nodes (2,157 against 2,191 on the first 30-site instance) and by a few seconds.
The root gap, not the tree code, is what the McCormick master cannot overcome.

\paragraph{Lazy McCormick rows.} The McCormick rows of the linked master (Section~\ref{sec:linked}) can be added only
when violated, like the lifted cuts. In the C++ code on the 50-site instance of seed 2 this takes the run from 277 s
and 87 nodes to 226 s and 65 nodes with the same optimum and the same final row count ($\approx$3,800); on the sparse
40-site instance below the row count falls from 986 to 899 at equal time. The rows are cheap and mostly tight, so the
gain is modest.

\paragraph{Where the time goes.} In the decomposed solver the master LP solve is 95--100\% of the wall time on the
40- and 50-site instances of Table~\ref{tab:main} (0.07--0.10 s per solve at $n = 40$, 0.15--0.20 s at $n = 50$,
12--24 solves per node), the separation with eight threads 0--5\%, and the closed-form block evaluation is not
measurable against them: the master LP is the cost.

\begin{table}[ht]
\centering\small
\caption{Sensitivity to the number of scenarios and to the endogeneity, $n = 30$, $m = 15$: decomposed solver (C++,
two threads) and node-local \textsc{MC-B} (C++, $S$ recourse LPs per node, wall 3,600 s); the $S = 20$ strong rows are the instances of
Table~\ref{tab:main}. ``Weak'' endogeneity: uniform base probabilities, each opened site moving at most $0.9/(nS)$
of mass per scenario along a linear ramp across scenarios.}
\label{tab:sens}
\begin{tabular}{llr rrr rrr}
\toprule
& & & \multicolumn{3}{c}{decomposed} & \multicolumn{3}{c}{\textsc{MC-B}, node-local}\\
\cmidrule(lr){4-6}\cmidrule(lr){7-9}
endogeneity & $S$ & seed & root gap & nodes & time (s) & root gap & nodes & time (s)\\
\midrule
strong & 20  & 1 & 0.29\% & 63 & 36 & 6.75\% & 2,191 & 101\\
strong & 20  & 2 & 0.07\% & 55 & 33 & 7.76\% & 663 & 15\\
strong & 20  & 3 & 0.22\% & 25 & 13 & 6.71\% & 1,531 & 39\\
strong & 50  & 1 & 0.26\% & 51 & 65 & 7.89\% & 1,239 & 164\\
strong & 50  & 2 & 0.26\% & 83 & 56 & 8.77\% & 2,535 & 608\\
strong & 50  & 3 & 0.24\% & 35 & 41 & 8.04\% & 1,599 & 239\\
strong & 100 & 1 & 0.25\% & 159 & 196 & 8.41\% & 4,270 & unproven (0.20\%)\\
strong & 100 & 2 & 0.18\% & 11 & 35 & 9.22\% & 3,189 & 2,235\\
strong & 100 & 3 & 0.30\% & 37 & 65 & 7.49\% & 2,111 & 1,351\\
weak   & 20  & 1 & 0.18\% & 141 & 22 & 0.79\% & 425 & 5\\
weak   & 20  & 2 & 0.00\% & 0 & 0 & 0.00\% & 1 & 0\\
weak   & 20  & 3 & 0.19\% & 123 & 43 & 0.71\% & 1,317 & 37\\
\bottomrule
\end{tabular}
\end{table}

\paragraph{Scenarios and endogeneity.} Table~\ref{tab:sens} varies what the generator holds fixed elsewhere. The root
gap is flat in $S$ (0.18--0.30\% at $S = 50$ and $100$ against 0.07--0.29\% at $S = 20$) and the time grows about
linearly in $S$ (means 27, 54, 99 s), as the per-scenario block structure of Section~\ref{sec:decomp} predicts: the
master gains $S$ epigraph variables and the separation $S$ recourse LPs, nothing quadratic. Weak endogeneity is
easier for the bound, not harder: with little mass moving, the products $p_s(y)Q_s(y)$ are nearly linear and the
root gap is at or below 0.2\% and even the McCormick gap falls to 0.8\%; on seed 2 both roots are integral. The
strong instances used throughout are the hard end of the generator. For the node-local McCormick master the
scenario count is what hurts: its root gap stays at 7.5--9.2\% and its time grows from 15--101 s at $S = 20$ to
1,351 s or more at $S = 100$, one instance unproven in an hour (gap 0.20\%), against 35--196 s for the decomposed
solver.

\begin{table}[ht]
\centering\small
\caption{Instances with $b_s = 0$ on the high-demand scenarios (generator option \texttt{--bpos0}; $m = 15$, $S = 20$,
strong endogeneity otherwise unchanged), on which the linking rows of those scenarios are vacuous. Root gaps
$(F^* - v)/F^*$ of the three masters at the converged root LP, and the trees of the decomposed solver and of node-local
\textsc{MC-B} (wall 3,600 s; the McCormick tree is run at 25 and 30 sites). $^\dagger$Gap against the incumbent of an
unproven run, an upper bound on the true gap.}
\label{tab:bpos0}
\begin{tabular}{ll rrr rr rr}
\toprule
& & \multicolumn{3}{c}{root gap} & \multicolumn{2}{c}{decomposed} & \multicolumn{2}{c}{\textsc{MC-B}, node-local}\\
\cmidrule(lr){3-5}\cmidrule(lr){6-7}\cmidrule(lr){8-9}
$n$ & seed & \textsc{MC-B} & unlinked & linked & nodes & time (s) & nodes & time (s)\\
\midrule
% TAB4-BEGIN
25 & 1 & 23.95\% & 1.26\% & 1.20\% & 259 & 56 & 715 & 16\\
25 & 2 & 19.15\% & 1.17\% & 1.13\% & 83 & 21 & 2,235 & 105\\
25 & 3 & 23.87\% & 1.14\% & 1.05\% & 77 & 63 & 1,273 & 32\\
30 & 1 & 26.25\% & 1.79\% & 1.71\% & 237 & 595 & 1,697 & 79\\
30 & 2 & 32.26\% & 3.16\% & 3.05\% & 275 & 91 & 2,223 & 156\\
30 & 3 & 16.94\% & 0.63\% & 0.57\% & 273 & 122 & 9,738 & unproven (0.74\%)\\
35 & 1 & 22.83\% & 1.39\% & 1.32\% & 621 & 2,818 &  & \\
35 & 2 & 27.44\% & 1.91\% & 1.83\% & 113 & 168 &  & \\
35 & 3 & 25.19\%$^\dagger$ & 1.77\%$^\dagger$ & 1.70\%$^\dagger$ & 537 & unproven (0.23\%) &  & \\
40 & 1 & 18.58\% & 0.83\% & 0.80\% & 771 & 2,379 &  & \\
40 & 2 & 17.19\% & 0.34\% & 0.29\% & 135 & 199 &  & \\
40 & 3 & 30.37\% & 1.44\% & 1.40\% & 415 & 2,834 &  & \\
% TAB4-END
\bottomrule
\end{tabular}
\end{table}

\paragraph{When the linking rows are vacuous.} The linking rows tighten scenario $s$ by exactly $\bpos_s(e^{0*}_s - L_0)$
(Remark~\ref{rem:bpos}), so a scenario with $\bpos_s = 0$ gains nothing from them. Table~\ref{tab:bpos0} takes this to the
limit: the generator variant \texttt{--bpos0} gives the high-demand half of the scenarios no base mass at all, so their
whole probability is moved in by opened sites and their linking rows are vacuous. Three things happen. The linked and
unlinked roots nearly coincide, within 0.1 points on every instance, as the remark predicts. The factor cuts alone still
hold the gap at 0.3--3.2\%, against 17--32\% for the McCormick master, whose global bounds are worst exactly when the
probability of a scenario can swing from zero. And the gap is larger than on the standard instances (0.07--0.43\%), so
the decomposed tree is larger and slower: 77--771 nodes and 21 s to 47 min, one 35-site instance unproven in an hour at
0.23\%. On a gap of this size the cheaper McCormick node pays: at 25 and 30 sites the node-local \textsc{MC-B} tree is
faster on four of six instances (16--156 s against 21--595 s) and unproven after an hour on the sixth (9,738 nodes,
0.74\%) where the decomposed solver needs 122 s; the McCormick tree was not run at 35 and 40 sites. These instances
mark the edge of the method: its advantage is the root, and when the root gap is in the percent range rather than the
tenth of a percent, the per-node master LP is expensive for what it closes at small $n$.

\begin{table}[ht]
\centering\small
\caption{Sensitivity to the number of customers, $n = 30$, $S = 20$, strong endogeneity, fresh instances from the
C++ generator, decomposed solver with eight threads, wall 3,600 s. ``Pool'' is the number of dual vertices in the pool
at the end of the root loop, ``added'' the vertices the tree adds to it. $^\dagger$Gap against the incumbent of an
unproven run, an upper bound on the true gap.}
\label{tab:msens}
\begin{tabular}{ll rrrr rr}
\toprule
$m$ & seed & root gap & root (s) & pool & added & nodes & time (s)\\
\midrule
% TAB5-BEGIN
15 & 1 & 0.30\% & 10 & 8,041 & 56 & 111 & 61\\
15 & 2 & 0.00\% & 3 & 2,196 & 0 & 3 & 3\\
15 & 3 & 0.31\% & 7 & 3,791 & 329 & 233 & 92\\
50 & 1 & 0.50\% & 24 & 16,008 & 1,018 & 985 & 325\\
50 & 2 & 0.57\% & 58 & 19,635 & 456 & 131 & 142\\
50 & 3 & 0.78\% & 58 & 16,806 & 2,083 & 397 & 282\\
100 & 1 & 1.29\% & 167 & 31,585 & 23,328 & 3,485 & 1,950\\
100 & 2 & 0.44\% & 77 & 29,385 & 202 & 63 & 115\\
100 & 3 & 0.83\% & 58 & 20,716 & 4,293 & 781 & 382\\
200 & 1 & 1.24\%$^\dagger$ & 203 & 42,640 & 53,934 & 3,858 & unproven (0.18\%)\\
200 & 2 & 1.09\% & 285 & 51,539 & 9,433 & 1,431 & 1,113\\
200 & 3 & 1.52\%$^\dagger$ & 272 & 48,914 & 45,334 & 5,302 & unproven (0.02\%)\\
% TAB5-END
\bottomrule
\end{tabular}
\end{table}

\paragraph{Customers.} Table~\ref{tab:msens} varies the number of customers at 30 sites on fresh instances from the
C++ generator (the $m = 15$ rows are therefore not those of Table~\ref{tab:main}). Each customer adds a row to every
recourse LP and a coordinate to every dual vertex, and it is the vertex count that grows: the root pool holds 2--8
thousand vertices at $m = 15$ and 43--52 thousand at $m = 200$, the root loop takes 3--10 s against 203--285 s, and the
separation itself, $S(1+2n)$ recourse LPs of size $n \times m$ per round, stays below 6\% of the time at every $m$. The
root gap grows with $m$ as well, from at most 0.31\% at $m = 15$ to 1.1--1.5\% at $m = 200$: a recourse function with
more pieces is less well approximated at the fractional root by cuts that are affine in the lifted variables. The
consequence is in the tree. Up to 100 customers every instance is proven (3 to 3,485 nodes, 3 s to 33 min); at 200
one instance is proven in 19 min and two are unproven after an hour at 0.18\% and 0.02\%, with 45--54 thousand vertices
added to the pool in the tree. The time goes where it always does: on the two unproven instances 3,446 and 3,397 of
the 3,600 s are master LP solves, 20--30 per node, since the node cutting-plane loop needs more iterations on a larger
gap. The number of customers reaches the solver through the gap and the pool, not through the size of the recourse LPs.

\section{Scaling and parallelism}\label{sec:scaling}

The dense lifting is quadratic in $n$: every cut row is about 60\% dense over the $n(n-1)/2$ pair columns, so master
nonzeros grow with $n^2$ and the simplex iteration cost with them (Table~\ref{tab:scale}). That is the one cost a
larger machine does not remove, and it sets what the dense master can do on a laptop.

\begin{table}[ht]
\centering\small
\caption{Dense master size at 4,000 rows; measured at $n = 50$, extrapolated on the 30--50-site trend. At $n = 100$
the measured root takes 15,200 rows and 11--21 minutes.}
\label{tab:scale}
\begin{tabular}{rrrrrr}
\toprule
sites & pairs & nnz/cut row & master nnz & memory/rank & s/solve\\
\midrule
50 (measured) & 1,225 & 720 & 2.9 M & 40 MB & 0.15--0.20\\
100 & 4,950 & 2,900 & 12 M & 150 MB & 1--3\\
200 & 19,900 & 11,700 & 47 M & 600 MB & 10--25\\
500 & 124,750 & 73,000 & 290 M & 3.5 GB & 60--250\\
\bottomrule
\end{tabular}
\end{table}

\paragraph{Node-parallel tree.} The tree parallelizes well: nodes are independent given the shared pool, and the few
cuts a node generates (only at integer leaves and rounded points) are valid everywhere. The node-parallel
implementation is an MPI program in which rank 0 computes the root once with all eight separation threads, then
broadcasts the dual-vertex pool, the active master cut rows and the stabilization center (two megabytes at 50 sites)
so that every worker starts from the converged root state; rank 0 then keeps the best-first heap, the incumbent and
the global pool, and each other rank owns a master LP and a replica of the pool, receives a node box together with any
cuts it has not yet seen, bounds it with the machinery of Section~\ref{sec:decomp}, and returns the branching
variable, any incumbent and any new cuts. Table~\ref{tab:mpi} gives the measurements with one coordinator and seven
workers on the otherwise idle 8-core laptop, against the single-rank code using all eight cores for separation. The
root costs the same in both (the idle ranks must be told to yield the processor while they wait for the broadcast,
otherwise their polling slows the root by half); the tree speedup on seven workers is $2.1$--$4.0\times$, lowest on the
trees of 25--91 nodes, where best-first search rarely has seven nodes of comparable bound open at once, and
$3.2$--$4.0\times$ on the 105--269-node trees. The net effect is $1.6$--$3.3\times$ end to end, growing with the tree,
which is the regime where parallelism is needed.

\begin{table}[ht]
\centering\small
\caption{MPI node-parallel tree, one coordinator and seven workers on an idle 8-core machine, root computed once on
the coordinator, against the single-rank code with eight separation threads. Same optima, all proven; ``tree'' is
the wall time after the root.}
\label{tab:mpi}
\begin{tabular}{lr rrr rrr rr}
\toprule
& & \multicolumn{3}{c}{single rank} & \multicolumn{3}{c}{$1{+}7$ ranks} & \multicolumn{2}{c}{speedup}\\
\cmidrule(lr){3-5}\cmidrule(lr){6-8}\cmidrule(lr){9-10}
$n$ & seed & nodes & root (s) & tree (s) & nodes & root (s) & tree (s) & tree & total\\
\midrule
40 & 1 & 21 & 17 & 36 & 25 & 16 & 17 & $2.1\times$ & $1.6\times$\\
40 & 2 & 91 & 18 & 76 & 91 & 20 & 36 & $2.1\times$ & $1.7\times$\\
40 & 3 & 101 & 12 & 143 & 105 & 12 & 37 & $3.9\times$ & $3.2\times$\\
50 & 1 & 125 & 35 & 359 & 113 & 29 & 92 & $3.9\times$ & $3.3\times$\\
50 & 2 & 63 & 22 & 171 & 63 & 24 & 54 & $3.2\times$ & $2.5\times$\\
50 & 3 & 281 & 38 & 668 & 269 & 43 & 169 & $4.0\times$ & $3.3\times$\\
\bottomrule
\end{tabular}
\end{table}

\paragraph{Dense 100 sites.} The C++ root on the three dense 100-site instances closes in 1,185, 680 and 1,277 s at
15,200--15,300 master rows, with root gaps to the best known solutions of 0.08\%, 0.04\% and 0.11\%: the gap does not
widen with $n$ at fixed $S$ and $m$. The tree is another matter. On seed 1 with one coordinator and seven workers the
search reached 100 nodes after 79 minutes with 89 nodes still open and a gap of 0.058\% (bounds 71.3123 and 71.3536),
and was stopped. At a node gap that small many boxes tie, the best-first heap widens, and each node still costs a
15,000-row master solve; closing a dense 100-site instance is a multi-hour job on eight cores and we report it as the
practical limit of the dense master on this hardware, not as a solved row.

\paragraph{Sparse influence at 100--500 sites.} At 500 sites with every site influencing demand, one facility shifts
0.2\% of the probability mass, which is not a model anyone would build; the realistic structure is \emph{sparse
influence}, a few hubs that change the demand regime. The sparse master of Section~\ref{sec:linked} keeps only
(influencing)$\times$(any) pair columns. Table~\ref{tab:sparse} uses the same generator with an influence density
$\rho$ chosen so that about ten sites influence demand at each size ($\rho = 0.1$, $0.05$, $0.02$); the sparse
master then has $O(|\mathcal I|n)$ pair columns instead of $O(n^2)$ and its $S$ blocks cost $O(S|\mathcal I|n)$ to
evaluate instead of $O(Sn^2)$, and 500 sites is a 26--222 s job on the
laptop with a root gap below 0.005\%. The dense and sparse masters give the same optimum where both were run (the
100-site instance of seed 1: dense root 28 s and 14,983 rows at a 0.07\% gap, sparse root 1 s and 2,164 rows at
0.08\%). The last row raises the number of customers from 15 to 100: the root gap is still below 0.005\% and
the root takes 108 s, but each node then costs about 1.5 s of recourse LPs and the tree is wide (2,865 nodes in two
hours, gap 0.003\% at the wall), so the recourse, not the master, becomes the cost. \textsc{MC-B} has only $nS$
bilinear terms and stays small at any size, so it remains the comparison at every scale; its root gap is unchanged by
sparsity.

\begin{table}[ht]
\centering\footnotesize
\caption{Sparse-influence instances, $m = 15$, $S = 20$, about ten influencing sites per instance; decomposed C++
solver with the sparse master, eight threads, wall 7,200 s.}
\label{tab:sparse}
\begin{tabular}{rlr rrrrr}
\toprule
$n$ & $\rho$ & seed & optimum & root gap (s) & nodes & time (s) & master rows\\
\midrule
100 & 0.10 & 1 & 47.829147 & 0.08\% (1) & 255 & 8 & 2,164\\
100 & 0.10 & 2 & 46.275180 & 0.04\% (1) & 121 & 4 & 2,184\\
100 & 0.10 & 3 & 50.156085 & 0.04\% (1) & 43 & 6 & 3,610\\
200 & 0.05 & 1 & 46.184683 & 0.01\% (3) & 35 & 8 & 6,112\\
200 & 0.05 & 2 & 46.601454 & 0.01\% (3) & 2,155 & 114 & 4,901\\
200 & 0.05 & 3 & 47.438137 & 0.01\% (3) & 667 & 59 & 4,448\\
500 & 0.02 & 1 & 47.766355 & $<$0.005\% (12) & 57 & 26 & 12,207\\
500 & 0.02 & 2 & 46.519291 & $<$0.005\% (31) & 385 & 222 & 20,922\\
500 & 0.02 & 3 & 45.413226 & $<$0.005\% (17) & 167 & 82 & 19,723\\
500, $m{=}100$ & 0.02 & 1 & 262.941205 & $<$0.005\% (108) & 2,865 & unproven (0.003\%) & 19,788\\
\bottomrule
\end{tabular}
\end{table}

\section{Related work}\label{sec:related}

Dupa\v{c}ov\'a's classification separates exogenous from endogenous uncertainty and, within the latter, the case of a
fixed scenario set with decision-dependent weights; she notes that convexity in the decision is lost and points to
enumeration and branch-and-bound \cite{dupacova}. Hellemo, Barton and Tomasgard formulate exactly the scenario-weight
model \eqref{eq:F} with linear-scaling and convex-combination forms of $p_s(y)$ and solve the bilinear programs with a
global solver \cite{hellemo}. Ahmed casts discrete first-stage problems with decision-dependent uncertainty as 0--1
hyperbolic programs \cite{ahmed}. Basciftci, Ahmed and Shen linearize binary-times-bounded-continuous products with
the McCormick set and note its exactness on binary points \cite{basciftci}; that is \textsc{MC-B}. Pantuso's
facility-location model with combinatorially many decision-dependent distributions enumerates distribution patterns
and linearizes an indicator times the second-stage cost \cite{fl2025}, and Pantuso and Hewitt's L-shaped method for
decision-dependent uncertainty uses McCormick envelopes for a distribution-independent LP bound \cite{pantuso}. In terms of scale, the exact methods above report at most 10 candidate sites in a monolithic McCormick MILP
\cite{basciftci} and, for the pattern-enumeration model, 20--25 candidate sites with 10--15 zones that select the
distribution ($2^{10}$--$2^{15}$ patterns) in minutes to two hours, with 17 zones solved partly and 20 not at all in
ten hours \cite{fl2025}. The models differ from \eqref{eq:F}, so Section~\ref{sec:results} is a comparison of scale
and structure, not of methods on common instances: there every one of 25--50 sites influences the probabilities, and
the sparse 500-site instances have about ten influencing sites. The lifting of Section~\ref{sec:linked} is the mixed-binary
reformulation--linearization of Adams and Sherali \cite{adams-sherali} on the dual side of the recourse; the
in-out stabilization is that of Ben-Ameur and Neto \cite{benameur}.



\section{Conclusion}

A decision-dependent probability does not have to be relaxed through global bounds on the recourse. Written as a
posiform, each of its terms multiplies a recourse LP by a binary variable, which is again a recourse LP at a scaled
right-hand side, and the shared dual polyhedron turns every dual vertex into a cut that is linear in the lifted
variables. The resulting LP is nearly tight at the root, decomposes by scenario with closed-form blocks, and its cuts
survive the whole branch-and-bound tree. On 25 instances with 25 to 50 sites it proves optimality in minutes
where the McCormick master in the same tree code does not finish in hours, and $1.6$--$13\times$ faster than the
McCormick master inside a full MIP solver. The fair McCormick
baseline, with recourse bounds recomputed at every node, is competitive up to 35 sites and then falls behind by a
factor that doubles every five sites, because its root gap is unchanged. The limit of the method is the same quantity
seen from the other side: on instances where the linking rows are vacuous ($b_s = 0$) or with 200 customers, the root
gap is in the percent range, the tree grows to thousands of nodes, and at 25--30 sites the cheaper McCormick node is
competitive again (Tables~\ref{tab:bpos0} and~\ref{tab:msens}).

The affine probability model is a modelling choice with limits: the influence of the sites is additive and must
remain a probability at every configuration, so it cannot express saturation or thresholds, which the conditional
distributions of \cite{fl2025} can. The construction does not depend on affinity, only on $p_s$ being a polynomial
in $y$: a multilinear $p_s$, which is every function of binary $y$, has a posiform in products of literals, each
product times $Q_s(y)$ is again the recourse LP at a right-hand side affine in the corresponding higher-order lifted
variable, and the same dual vertices give the cuts, at the price of level-$k$ RLT lifting. We have not pursued this.
The quadratic lifting bounds dense instances
near 100 sites on a laptop; with sparse influence, the realistic structure at that scale, the same solver proves
500-site instances in minutes.

\section*{Declaration on the use of AI assistance}

This work was carried out with Claude (Anthropic), used as an AI assistant throughout the project: for the
cross-field literature survey from which the lifted-cut construction was selected, for writing the C++ solver and
the experiment scripts, for running and recording the computational study, and for drafting and revising this
text, including the responses to review. The research question, the choice of method, the standards for the
experiments and the final decisions on every claim are the author's. The author reviewed and edited all generated
content, verified the proofs and the reported numbers against the solver's output, and takes full responsibility for
the content of the publication. The source code, the instances and the result files that support every number in
the paper are provided with it.

\begin{thebibliography}{9}
\bibitem{adams-sherali} W.~P.~Adams and H.~D.~Sherali. Linearization strategies for a class of zero-one mixed
integer programming problems. \emph{Operations Research} 38(2):217--226, 1990.
\bibitem{ahmed} S.~Ahmed. \emph{Strategic planning under uncertainty: stochastic integer programming approaches.}
PhD thesis, University of Illinois at Urbana-Champaign, 2000.
\bibitem{basciftci} B.~Basciftci, S.~Ahmed and S.~Shen. Distributionally robust facility location problem under
decision-dependent stochastic demand. \emph{European Journal of Operational Research} 292(2):548--561, 2021.
\bibitem{benameur} W.~Ben-Ameur and J.~Neto. Acceleration of cutting-plane and column generation algorithms:
applications to network design. \emph{Networks} 49(1):3--17, 2007.
\bibitem{dupacova} J.~Dupa\v{c}ov\'a. Optimization under exogenous and endogenous uncertainty. In \emph{Proceedings
of MME 2006}. \url{https://www2.karlin.mff.cuni.cz/~kopa/papers/MME06sty.pdf}
\bibitem{fl2025} G.~Pantuso. Solution of stochastic facility location problems with combinatorially many
decision-dependent distributions. \emph{Journal of Combinatorial Optimization}, 2025. doi:10.1007/s10878-025-01371-7.
\bibitem{hellemo} L.~Hellemo, P.~I.~Barton and A.~Tomasgard. Decision-dependent probabilities in stochastic
programs with recourse. \emph{Computational Management Science} 15:369--395, 2018.
\bibitem{highs} Q.~Huangfu and J.~A.~J.~Hall. Parallelizing the dual revised simplex method. \emph{Mathematical
Programming Computation} 10:119--142, 2018.
\bibitem{pantuso} G.~Pantuso and M.~Hewitt. The L-shaped method for stochastic programs with decision-dependent
uncertainty. \emph{Mathematical Programming}, 2025. doi:10.1007/s10107-025-02246-9.
\end{thebibliography}

\end{document}
